Large language models (LLMs) have demonstrated remarkable capabilities across diverse tasks, yet their deployment faces three critical challenges: unreliable outputs that may contain errors or hallucinations, limited control over generation behavior, and inefficient training procedures that waste computational resources on redundant data. This thesis addresses these challenges through a unified approach that leverages signals from the model's forward pass---both internal representations and output distributions---to improve reliability, control, and training efficiency.

We present three interconnected studies. First, we investigate uncertainty estimation by analyzing the relationship between token-level entropy and answer correctness. Our experiments on MMLU-Pro across four models (Phi-4, Mistral-Small-24B, Qwen-1.5B, Qwen-72B) demonstrate that single-token answer entropy achieves ROC-AUC scores up to 0.83 for predicting incorrect answers in knowledge-dependent domains, while Model-as-Judge approaches perform near-randomly. Second, we explore the geometric structure of multilingual representations in LLM latent space, discovering that language-specific directions can be identified through PCA with 96--100\% classification accuracy across four language pairs. By projecting out the source language component from hidden states, we reduce distributional divergence by up to 55\% and achieve 63--95\% reduction in Code-Switching Index in generated text, without requiring additional training. Third, we develop a complexity-aware fine-tuning pipeline that uses entropy to stratify training data and applies targeted interventions: direct supervised fine-tuning for regular examples and chain-of-thought distillation for complex ones. This approach achieves accuracy improvements from 0.39 to 0.52 (Qwen-3B) and 0.51 to 0.64 (Phi-4-Mini) on MMLU-Pro while using 79--82\% less training data, with results generalizing to GSM8K and MedMCQA datasets.

Together, these contributions demonstrate that accessible signals from a single forward pass provide valuable information for addressing fundamental challenges in LLM deployment. The methods developed require no architectural modifications or expensive retraining, making them practical for real-world applications.
